\documentclass[11pt]{article}
\usepackage[a4paper, margin=1in]{geometry}
\usepackage{amsmath, amssymb}
\usepackage{hyperref}
\usepackage{xcolor}
\usepackage{enumitem}
\usepackage{titlesec}
\usepackage{booktabs}
\usepackage{tabularx}
\hypersetup{colorlinks=true, linkcolor=blue!50!black, urlcolor=blue!50!black}
\titleformat{\section}{\large\bfseries\color{blue!40!black}}{\thesection.}{0.5em}{}
\titleformat{\subsection}{\normalsize\bfseries}{\thesubsection}{0.5em}{}

\title{\bfseries Interpreting Results: From Numbers to Real-World Contribution\\[0.3em]
\large How to read what your analyses say, and how to argue that the work matters}
\author{}
\date{}

\begin{document}
\maketitle

\section*{The interpretive task}
You will produce a lot of numbers. Mathematicians are tempted to report them all, leaving the reader to draw conclusions. \emph{This is a publication-killing habit in applied ML.} The reader needs you to translate numbers into claims, and claims into clinical or scientific contribution. This document gives you the framework for doing that.

\section{The five layers of interpretation}

For any quantitative result, ask the following five questions in order:

\begin{enumerate}
\item \textbf{What does the metric say literally?} (e.g.\ NRMSE = 0.05)
\item \textbf{Is the metric appropriate?} (would a clinician care about this metric, or about something else like peak-WSS error?)
\item \textbf{What does this metric translate to physically?} (5\% RMS error in WSS in absolute Pa terms; comparison to inter-patient variability)
\item \textbf{What does this enable, clinically or scientifically?} (acceptable for screening; not yet acceptable for surgical planning)
\item \textbf{What are the caveats?} (this regime; this geometry; this noise model)
\end{enumerate}

\textbf{Every interpretation in your paper should walk these five rungs.} If you stop at rung 1, the paper feels disconnected from the world. If you skip to rung 4, you've over-claimed.

\section{Interpretation by metric}

\subsection*{NRMSE on velocity}
\textbf{Literal:} the RMS error normalised by the RMS of the truth, often expressed as a percentage.\\
\textbf{Appropriate?} For the velocity field as a whole, yes. For WSS specifically, no --- WSS depends on derivatives, so a low velocity NRMSE does not imply a low WSS error.\\
\textbf{Physical translation:} NRMSE 5\% on a peak velocity of 1\,m/s is roughly 5\,cm/s error --- comparable to the noise floor of well-acquired clinical 4D-flow.\\
\textbf{What it enables:} cohort-level summary metrics; flow-rate computation.\\
\textbf{Caveat:} \emph{always} report a derivative-sensitive companion metric (WSS error or near-wall NRMSE).

\subsection*{$R^2$ coefficient of determination}
\textbf{Literal:} fraction of variance explained.\\
\textbf{Appropriate?} Useful as a one-number summary but \emph{misleading at high values}. The source paper reports $R^2 \geq 0.997$ on velocity components but still has WSS errors. $R^2$ near 1 means the gross structure is right; it says little about steep gradients.\\
\textbf{What it enables:} reassurance that nothing is catastrophically wrong.\\
\textbf{Caveat:} treat as a sanity check, never as a headline.

\subsection*{Pointwise relative error in WSS}
\textbf{Literal:} $|\tau_w^{\text{pred}} - \tau_w^{\text{true}}|/\tau_w^{\text{true}}$ at each wall point.\\
\textbf{Appropriate?} Yes --- this is the real headline.\\
\textbf{Physical translation:} 20\% relative error in mean aortic TAWSS ($\sim 1$\,Pa) is about 0.2\,Pa, which is roughly the inter-patient natural variability. So 20\% error is at the edge of clinical usefulness; 50\% is not.\\
\textbf{What it enables:} ranked: $<10\%$ enables risk-stratification research; 10--20\% enables population-level studies; 20--40\% enables only qualitative pattern recognition; $>40\%$ is essentially noise.\\
\textbf{Caveat:} relative error blows up where $\tau_w \to 0$ (recirculation, stagnation). Always also report the absolute error in Pa.

\subsection*{Peak WSS error}
\textbf{Literal:} the difference in maximum WSS over the surface.\\
\textbf{Appropriate?} Critically yes for stenoses (peak-WSS region is the rupture-relevant location); critically yes for aneurysms (peak-WSS predicts wall remodelling).\\
\textbf{Physical translation:} a peak-WSS error of 30\% understates plaque-rupture risk by an unknown but physiologically important amount.\\
\textbf{What it enables:} a low peak-WSS error is a strong specific claim and reads well in a clinical bridge journal.\\
\textbf{Caveat:} peak is a fragile statistic. Report also a high quantile (e.g.\ 95th percentile) which is more stable.

\subsection*{Mean WSS error / TAWSS error}
\textbf{Literal:} aggregate over wall surface.\\
\textbf{Appropriate?} Yes for atherogenesis-related applications (low TAWSS regions matter).\\
\textbf{Physical translation:} 0.1\,Pa is comparable to clinical 4D-flow precision; 0.5\,Pa is well above natural noise.\\
\textbf{What it enables:} statements about the magnitude of clinical bias.\\
\textbf{Caveat:} averaging hides regional pathology. Always pair with a regional analysis (Bland--Altman, surface error map).

\subsection*{Bland--Altman bias and limits of agreement}
\textbf{Literal:} systematic offset (bias) and 95\% spread (limits of agreement, LoA).\\
\textbf{Appropriate?} The most clinically literate scalar pair you can report.\\
\textbf{Physical translation:} a bias of $-0.15$\,Pa with LoA $[-0.6, 0.3]$ Pa on TAWSS means: ``on average we under-estimate by 0.15\,Pa, but in the worst 2.5\% of cases we under-estimate by 0.6 or over-estimate by 0.3.''\\
\textbf{What it enables:} direct comparability with the clinical 4D-flow-vs-CFD literature, where Bland--Altman is the lingua franca.\\
\textbf{Caveat:} requires sufficiently many comparison points (per-vertex; or per-region across geometries).

\subsection*{Conservation residuals (mass, momentum)}
\textbf{Literal:} the deviation of the trained network from $\nabla\cdot\mathbf{u} = 0$ or from cross-section flow-rate constancy.\\
\textbf{Appropriate?} Always report. Cheap to compute, hard to fake.\\
\textbf{Physical translation:} 1\% deviation = the network does not conserve mass perfectly but is well within noise. 10\% = the network is exploiting the loss function in non-physical ways and the velocity field is suspect.\\
\textbf{What it enables:} reviewer trust.\\
\textbf{Caveat:} a network can have small residuals and still be wrong; conservation is necessary, not sufficient.

\subsection*{Computational time}
\textbf{Literal:} CFD vs.\ PINN training vs.\ PINN inference, all in wall-clock seconds.\\
\textbf{Appropriate?} Yes, but the comparison must be \emph{fair}. CFD runs once per case; PINN trains once per case (or once per family). State the amortisation explicitly.\\
\textbf{Physical translation:} ``CFD takes 2 hours per case; PINN training 30 min per case; PINN inference 10 ms per query. For workflows with $>10$ queries per geometry, PINN inference dominates.''\\
\textbf{What it enables:} the ``why this matters in practice'' argument.\\
\textbf{Caveat:} training time is a real cost --- don't bury it in a footnote.

\section{Interpreting the central artefact: the (voxel size, SNR) phase diagram}

This is your central claim. The figure is a 2D heatmap; your job is to extract three or four headline numbers from it.

\subsection*{The reliable region}
Define a threshold (e.g.\ NRMSE on $\tau_w$ < 20\%, or peak-WSS error $< 25\%$) and report the largest voxel size at each SNR level inside the reliable region. The result is a curve in the (voxel, SNR) plane. \emph{This curve is the answer to ``how good must MRI be to make WSS recovery via PINN reliable?''} --- the question your paper is for.

\subsection*{Comparison to clinical 4D-flow specifications}
Position the typical clinical operating point on the diagram (voxel 1.5--2.5\,mm; VNR 5--15). Is it inside or outside the reliable region? Honest answers in either direction make a paper.

\textbf{If outside:} ``current clinical 4D-flow does not reliably support PINN-based WSS recovery for this geometry class; we identify the resolution/SNR target needed.''\\
\textbf{If inside:} ``current clinical 4D-flow specifications are sufficient for reliable PINN-based WSS recovery in this geometry class.''\\
\textbf{Mixed:} ``the boundary cuts through the clinical operating range; we identify the geometry-specific minimum specification.''

\subsection*{Methodological-tweak shift}
With your tweak (RFF, hard no-slip, Bayesian) on, the reliable region expands or shifts. Quantify: ``the tweak buys $\Delta$\,mm of voxel size at SNR = 10'' or ``at clinical operating point X, the WSS error drops from Y\% to Z\%.''

\subsection*{Geometry sensitivity}
Compare the boundary across the three geometries. Is concentric stenosis easier than eccentric? Are bifurcations harder than straight stenoses? \emph{The geometry-by-geometry differences are themselves a contribution.}

\section{Interpreting Bayesian uncertainty (if you go this route)}

\subsection*{Calibration}
Plot empirical coverage vs.\ nominal credible level. If your model says ``80\% credible interval'' and 80\% of true values are inside it, you are calibrated. Deviations:
\begin{itemize}[noitemsep]
\item \textbf{Over-confident:} 80\% interval contains only 50\% --- the model under-estimates uncertainty. Common pathology; clinically dangerous.
\item \textbf{Under-confident:} 80\% interval contains 95\% --- conservative; safe but uninformative.
\item \textbf{Roughly calibrated:} on diagonal $\pm$10\%. Acceptable.
\end{itemize}

\subsection*{Spatial structure of uncertainty}
The uncertainty map should show \emph{high} variance in the regions where you'd expect failure: near the wall, in recirculation zones, at the throat. If it doesn't, your model is over-confident in problematic regions and you should say so.

\subsection*{Correlation with truth error}
Compute the correlation between predicted standard deviation and actual error magnitude over wall points. A high correlation ($\rho \gtrsim 0.5$) is the strong claim: ``where the model says it's uncertain, it actually is wrong.''

\section{Demonstrating real-world contribution}

The paper succeeds if you can answer, in a single paragraph, the question: \emph{what could a clinician or applied scientist do tomorrow that they could not do today, because of this paper?} Your candidate answers, ranked by strength:

\subsection*{The strongest claim: an MRI specification recipe}
``A 4D-flow MRI acquisition with voxel $\leq 1.5$\,mm and VNR $\geq 10$ on a stenosed carotid is sufficient to recover TAWSS to within 15\% of CFD; this enables retrospective TAWSS analysis of existing 4D-flow studies and provides design targets for new acquisitions.'' This is concrete, falsifiable, and useful.

\subsection*{The methods claim}
``We provide an open implementation of a PINN-based WSS estimator that matches or beats existing 4D-flow WSS extraction methods (PaLMA, divergence-free smoothing) on benchmark synthetic cases.'' This claim becomes stronger if your code is genuinely reusable.

\subsection*{The benchmark claim}
``We release a benchmark of CFD ground truth + synthetic 4D-flow at controlled operating points; future methods can be compared on a common task.'' This is how the broader community gets pulled in. Highly leverageable.

\subsection*{The mathematical claim (if you take Problem 1 from the strategy doc on board)}
``We show that Fourier-feature embedding shifts the reliable region in a quantifiable way; the bandwidth $\sigma$ should be tuned to the voxel scale.'' This is the closest thing to a theorem you'll get; rigorous justification (NTK-based or otherwise) makes it a real contribution.

\subsection*{The honest-limits claim}
``We characterise the regime where PINN-based WSS recovery is currently \emph{not} reliable, and identify three failure modes.'' This sounds modest but is unusually valuable because the field tends to oversell.

\section{Structuring the discussion section}

The Discussion is where the interpretation lives. Recommended structure:

\paragraph{Paragraph 1 --- Restate the principal findings.} Concrete numbers from the headline figures. Three sentences max. Tie to the headline claims C1--C4.

\paragraph{Paragraph 2 --- Comparison to prior work.} Position your operating-point characterisation against Arzani 2021, Fathi 2020, IP-PINN 2024, and the non-PINN WSS methods (PaLMA, PG-VGC). Be specific about what's new (the systematic sweep, the benchmark, the tweak ablation).

\paragraph{Paragraph 3 --- Clinical implications.} Translate the phase diagram into language a clinical 4D-flow user would recognise. Specify acquisition parameters. State which clinical questions this enables (e.g.\ atherogenesis-prone region screening, longitudinal aneurysm monitoring).

\paragraph{Paragraph 4 --- Methodological insight.} Why does the tweak help? Connect to the spectral-bias / NTK story. Mathematicians: this is the part you'll find natural.

\paragraph{Paragraph 5 --- Failure modes and limitations.} Be specific. Examples: ``WSS recovery fails in geometries with surface roughness amplitude $>X\%$ of vessel radius;'' ``the synthetic forward operator does not include displacement artefacts which may further bias real-MRI applications;'' ``rigid wall assumption excludes pulsatile-FSI effects.'' Naming limitations protects you from reviewer attack.

\paragraph{Paragraph 6 --- Future work.} Three concrete directions. Pulsatile (Womersley) extension, FSI, patient-specific cohort study. Don't list six --- list three you'd actually pursue.

\section{The cardinal sin of result interpretation in this field}

Reporting an extremely low NRMSE on velocity, claiming ``high accuracy,'' and never showing the WSS map. The source paper itself avoids this trap (it shows centreline and radial profiles), but many do not. Your reviewers will. Do not.

\section*{Self-check: before claiming contribution, ask}
\begin{itemize}[noitemsep]
\item Have I shown my method works (Figure 5: WSS maps)?
\item Have I shown when it works and when it doesn't (Figure 6: phase diagram)?
\item Have I quantified physical consistency (Figure 9: conservation)?
\item Have I shown calibrated uncertainty or admitted I haven't (Figure 11)?
\item Have I made a falsifiable claim about an acquisition specification?
\item Have I named at least three honest limitations?
\item Have I released code and data?
\end{itemize}
If yes to all, you have a paper. If no to any, fix that before submitting.

\end{document}
